\documentclass[11pt]{article}
\usepackage[margin=1in]{geometry}
\usepackage{amsmath,amssymb}
\usepackage{booktabs}
\usepackage{enumitem}

\newcommand{\R}{\mathcal{R}}   % reservoirs
\newcommand{\Cs}{\mathcal{C}}  % consumers
\newcommand{\N}{\mathcal{N}}   % nodes
\newcommand{\A}{\mathcal{A}}   % arcs

\title{A Water Distribution Network Design Model\\
\large An MPEC reformulation of the original MINLP}
\author{Mathematical description of \texttt{water-net.mod}}
\date{}

\begin{document}
\maketitle

\section{Overview}

The model designs a water distribution network: it chooses pipe diameters,
the supply drawn from reservoirs, and the resulting flows and pressures so as
to minimize the combined cost of pumping, raw water, and pipe investment, while
respecting flow conservation and the physical pressure-loss law on each pipe.

The original formulation is a mixed-integer nonlinear program (MINLP): a binary
variable on each arc selects the direction of flow. Here the flow on each arc is
split into a nonnegative \emph{positive} part $qp$ and a nonnegative
\emph{negative} part $qn$, and the integrality is replaced by a
\emph{complementarity} condition that forbids both parts from being active
simultaneously. The result is a continuous mathematical program with
complementarity constraints (MPEC), described below.

\section{Sets}

\begin{align*}
\N &= \text{nodes of the network},\\
\R &\subseteq \N : \text{reservoirs (supply sources)},\\
\Cs &\subseteq \N : \text{consumers (demand sinks)},\\
\A &\subseteq \N \times \N : \text{arcs (directed pipe connections)}.
\end{align*}

\section{Parameters}

Each node $i \in \N$ carries data
$\mathit{demand}_i$, $\mathit{height}_i$, coordinates $(x_i, y_i)$,
$\mathit{supply}_i$, $\mathit{wcost}_i$ (water cost), and
$\mathit{pcost}_i$ (pumping cost). The remaining parameters are scalars:

\begin{center}
\begin{tabular}{llr}
\toprule
Symbol & Meaning & Value \\
\midrule
$\mathit{dpow}$  & power on diameter in the pressure-loss equation & $5.33$ \\
$\mathit{dmin}$  & minimum pipe diameter & $0.15$ \\
$\mathit{dmax}$  & maximum pipe diameter & $2.00$ \\
$\mathit{hloss}$ & constant in the pressure-loss equation & $1.03\times10^{-3}$ \\
$\mathit{dprc}$  & scale factor in the investment-cost equation & $6.90\times10^{-2}$ \\
$\mathit{cpow}$  & power on diameter in the cost equation & $1.29$ \\
$r$              & interest rate & $0.10$ \\
$\mathit{maxq}$  & bound on $qp$ and $qn$ & $2.00$ \\
\bottomrule
\end{tabular}
\end{center}

Several parameters are derived from the data:
\begin{align}
\mathit{dist}_{ij} &= \sqrt{(x_i - x_j)^2 + (y_i - y_j)^2},
   && (i,j)\in\A \quad\text{(Euclidean pipe length)},\\
\mathit{davg} &= \sqrt{\mathit{dmin}\cdot\mathit{dmax}}
   && \text{(geometric-mean diameter, used for initialization)},\\
\mathit{rr} &= \frac{\sum_{i\in\Cs}\mathit{demand}_i}
                    {\sum_{i\in\R}\mathit{supply}_i}
   && \text{(ratio of total demand to total supply)},\\
\mathit{hl}_i &= \mathit{height}_i +
   \begin{cases} 7.5 + 5\,\mathit{demand}_i, & i \in \Cs,\\[2pt] 0, & i \notin \Cs,\end{cases}
   && i\in\N \quad\text{(minimum required pressure head)}.
\end{align}

\section{Variables}

\begin{align}
qp_{ij} &\in [0,\ \mathit{maxq}], && (i,j)\in\A &&\text{positive flow on each arc } [\mathrm{m^3/s}],\\
qn_{ij} &\in [0,\ \mathit{maxq}], && (i,j)\in\A &&\text{negative flow on each arc } [\mathrm{m^3/s}],\\
d_{ij}  &\in [\mathit{dmin},\ \mathit{dmax}], && (i,j)\in\A &&\text{pipe diameter } [\mathrm{m}],\\
h_i     &\ge \mathit{hl}_i, && i\in\N &&\text{pressure head at each node } [\mathrm{m}],\\
s_i     &\in [0,\ \mathit{supply}_i], && i\in\R &&\text{supply drawn at each reservoir } [\mathrm{m^3/s}].
\end{align}

\section{Objective}

The objective \texttt{cost} sums three terms: the discounted operating cost at
the reservoirs (pumping against the head $h_i-\mathit{height}_i$ plus raw-water
cost, annualized by the interest rate $r$), the pipe investment cost, and a
small penalty on total flow:
\begin{equation}
\tag{\texttt{cost}}
\min\quad
\frac{1}{r}\sum_{i\in\R} s_i\,\mathit{pcost}_i\,\bigl(h_i - \mathit{height}_i\bigr)
\;+\; \frac{1}{r}\sum_{i\in\R} s_i\,\mathit{wcost}_i
\;+\; \mathit{dprc}\sum_{(i,j)\in\A}\mathit{dist}_{ij}\,d_{ij}^{\,\mathit{cpow}}
\;+\; \sum_{(i,j)\in\A}\bigl(qp_{ij} + qn_{ij}\bigr).
\end{equation}

\section{Constraints}

\paragraph{Flow conservation (\texttt{cont}).}
At every node, net inflow plus any reservoir supply equals the demand:
\begin{equation}
\tag{\texttt{cont}}
\sum_{(j,i)\in\A}\bigl(qp_{ji} - qn_{ji}\bigr)
\;-\; \sum_{(i,j)\in\A}\bigl(qp_{ij} - qn_{ij}\bigr)
\;+\; \begin{cases} s_i, & i\in\R,\\ 0, & i\notin\R \end{cases}
\;=\; \mathit{demand}_i,
\qquad i\in\N.
\end{equation}
Here $\sum_{(j,i)\in\A}$ runs over arcs entering node $i$ and
$\sum_{(i,j)\in\A}$ over arcs leaving it; the signed flow on an arc is
$qp - qn$.

\paragraph{Pressure loss (\texttt{loss}).}
On each arc the head difference is related to the flow and diameter through the
nonlinear pressure-loss law (with flow exponent $\mathit{qpow}=2$):
\begin{equation}
\tag{\texttt{loss}}
h_i - h_j
\;=\; \mathit{hloss}\,\mathit{dist}_{ij}\,
   \frac{qp_{ij}^{\,2} - qn_{ij}^{\,2}}{d_{ij}^{\,\mathit{dpow}}},
\qquad (i,j)\in\A.
\end{equation}

\paragraph{Complementarity (\texttt{compl}).}
On each arc the flow is either positive or negative, never both. Since
$qp_{ij}, qn_{ij}\ge 0$, this is enforced by requiring their products to vanish,
written as a single aggregated inequality:
\begin{equation}
\tag{\texttt{compl}}
\sum_{(i,j)\in\A} qp_{ij}\,qn_{ij} \;\le\; 0.
\end{equation}
Together with nonnegativity this forces $qp_{ij}\,qn_{ij}=0$ for every arc,
i.e.\ $0 \le qp_{ij} \perp qn_{ij} \ge 0$, the complementarity that replaces the
binary direction variable of the original MINLP.

\section{Remarks}

The model is nonconvex: the pressure-loss constraint \texttt{loss} is bilinear
in the flows and divides by $d_{ij}^{\,5.33}$, and the investment term in
\texttt{cost} is a fractional power of the diameter. The complementarity
constraint \texttt{compl} makes the feasible region degenerate (its linear
independence constraint qualification fails at every feasible point), which is
the characteristic difficulty of MPECs.

\end{document}
